% chapters/ch00-example.tex -- FORMAT EXAMPLE, not a chapter of any book.
%
% Copy this file to chapters/chNN-<slug>.tex, replace everything in angle
% brackets, and put your own slug in place of SLUG in every label.  main.tex
% deliberately does NOT \include it: an unedited example must never reach the
% PDF.
%
% ASCII only under tex/.
\chapter{<Title: name the object, not the technique>}
\label{ch:SLUG}

% ---------------------------------------------------------------- header box
% Required of every chapter (CLAUDE.md, Phase 4).  Four items, in this order.
% Write all four in PLAIN WORDS: no depth grades (L0-L4), no node IDs, no
% baseline numbers, no Phase numbers, and no pointer to spec.md /
% background.md / progress.md.  Those are our working devices, not the
% reader's.  Where a grade used to justify an omission ("we stop here, it is
% L2"), write the reason instead ("the paper only quotes the statement, so we
% stop here").
\begin{chapterheader}
\textbf{Purpose.}
<What this chapter is for, in terms of the paper.  Name the objects it builds
and the one thing they are needed for.  Three to six sentences.  Not a table
of contents of the chapter.>

\medskip
\textbf{Assumed.}
<What the reader is taken to know already, listed concretely.  Written from
the knowledge diagnosis but never citing it -- and it is also the promise that
nothing K2-or-above gets re-explained below.  Say explicitly what is NOT used
here if that spares the reader a detour.>

\medskip
\textbf{Supports.}
<Which part of the paper this serves: section numbers, and what those sections
do with it.  If a later chapter depends on this one, say so with a
cross-reference.>

\medskip
\textbf{How far things are proved.}
<The honest scope.  "Everything, with complete proofs", or "statements and
their role only", or a mixture -- and the reason for the line where it falls.
If the paper imitates a proof rather than quoting a result, that is the reason
to prove it here, and it belongs in this paragraph.>
\end{chapterheader}

% A short orientation paragraph before the first section: conventions that hold
% throughout the chapter, and nothing else.
Throughout, <standing conventions: the letters, the standing hypotheses, the
one notational choice that differs from the reader's habit>.

%% =====================================================================
\section{<First section: the definition and what it is for>}
\label{sec:SLUG-first}

% Every numbered item carries an exact pointer.  Author and year is not enough
% (CLAUDE.md, verification rules).
%
% TRAP: inside the OPTIONAL argument of a theorem environment, \cite[...]{...}
% ends at the first ']'.  Brace the whole pointer, as below.
\begin{defn}[{\cite[Def.~<n.m>]{Har77}}]
\label{def:SLUG}
<The definition, in the paper's notation.  If the paper's notation departs
from the standard one, follow the paper and say so in a footnote.>
\end{defn}

<One or two sentences on why the definition is shaped this way -- the reader
should learn what the definition is guarding against.>

\begin{prop}[{\cite[Thm.~<n.m>]{Har77}}]
\label{prop:SLUG}
<The statement exactly as the paper uses it, not in the greatest generality
the source offers.>
\end{prop}

\begin{proof}
<Include a proof only when the chapter's stated scope calls for one.  If the
source proves it, say which source and reproduce the argument; if the source
leaves a gap that we fill, mark it.>
\OWNPROOF{the source states the result without proof; the argument below is
ours}
\end{proof}

\begin{ex}
\label{ex:SLUG}
<An example or a computation that the reader can check.  One per statement at
most, and only where the chapter's scope allows examples.>
\end{ex}

\begin{rem}
\label{rem:SLUG}
<A remark earns its place by explaining a choice, a failure of the obvious
generalisation, or the exact point where the paper will use this.  A remark
that only restates the definition is deleted.>
\end{rem}

%% =====================================================================
% The chapter closes with this table.  It is what makes the chapter auditable:
% every numbered item above must appear in it, and every row must name a place
% in the paper (or a later chapter) that uses the item.  An item with no
% "Used in" entry is either deleted or justified in the text as intuition the
% reader needs (CLAUDE.md, Phase 4, the check before a chapter is finalised).
\section{Summary: where this chapter is used}
\label{sec:SLUG-summary}

\begin{center}
\small
\begin{tabular}{@{}L{4.6cm}lL{5.9cm}@{}}
\toprule
Item & Here & Used in <PAPER> \\
\midrule
<the object or statement, named as the reader would look for it>
  & \ref{def:SLUG}
  & <\S of the paper; and in one clause, what it is used for there> \\
<the statement the paper actually applies>
  & \ref{prop:SLUG}
  & <\S of the paper; or a later chapter, by cross-reference> \\
<the worked example>
  & \ref{ex:SLUG}
  & <what it is the local model of, or which hypothesis it shows to be
    necessary> \\
\bottomrule
\end{tabular}
\end{center}

\noindent
<A closing paragraph of two or three sentences: what the chapter decided, and
what the reader should carry into the next one.  Optional, but it is the place
to say which single item of the table is the load-bearing one.>
